\documentclass[11pt]{article}
\usepackage[utf8]{inputenc}
\usepackage[T1]{fontenc}
\usepackage[margin=1in]{geometry}
\usepackage{hyperref}
\usepackage{enumitem}
\hypersetup{colorlinks=true,linkcolor=black,urlcolor=blue,citecolor=blue}
\title{LEGO Game and Conflict Resolution}
\author{Piter Garcia}
\date{March 20, 2026}
\begin{document}
\maketitle
\section*{Packet Use}
This printable lesson packet is generated from the paired Markdown guide. The Markdown version remains the clickable, neurodivergent-friendly working copy; this PDF carries the same lesson substance in a formal handout format. Evidence claims in this packet are based on transcript-derived lesson notes and the cleaned transcript files in this workspace.

\section*{Quick Scan}
\begin{itemize}[leftmargin=*]
\item \textbf{Grade:} Grade 1
\item \textbf{Grade Hub:} Notes [../grade-1-Notes.md]
\item \textbf{Schedule Reference:} Spring2026\_TeachingSchedule\_Derek.md [../../school/Spring2026\_TeachingSchedule\_Derek.md]
\item \textbf{Workbook Sheet:} \texttt{1st} in \texttt{Teaching\_Placement-IGNITE Curriculum\_ GCSD25-26.xlsx}
\item \textbf{Lesson Focus:} A structured build or game task used to teach teamwork, turn-taking, and reflective problem solving while students work with shared materials.
\item \textbf{CS Alignment Strength:} Inclusive computing culture / collaboration support lesson
\item \textbf{LaTeX Source:} lego-game-and-conflict-resolution.tex [./lego-game-and-conflict-resolution.tex]
\item \textbf{Bib File:} lego-game-and-conflict-resolution.bib [./lego-game-and-conflict-resolution.bib]
\item \textbf{Artifacts Folder:} artifacts/ [./artifacts/]
\end{itemize}

\section*{Lesson Identity}
\begin{itemize}[leftmargin=*]
\item \textbf{Likely Class Blocks:} Day 3 -- 12:25-1:15 -- Gargana, Day 3 -- 2:15-3:05 -- Polfleit, Day 5 -- 12:25-1:15 -- Montgomery
\item \textbf{Main Lesson Family:} LEGO Game and Conflict Resolution
\item \textbf{Primary Evidence Base:} Transcript-derived notes plus source transcripts from this teaching-placement workspace.
\item \textbf{Primary External Source Type:} Official curriculum/provider materials.
\end{itemize}

\section*{What We Are Teaching}
\begin{itemize}[leftmargin=*]
\item following shared task rules in a maker or game environment
\item using language to repair conflict during a collaborative build
\item breaking a challenge into steps and returning to the goal after a disruption
\item treating classroom materials as shared computing/maker resources
\end{itemize}

\section*{CS Standards Alignment}
\begin{itemize}[leftmargin=*]
\item \textbf{1A-IC-17}: Work respectfully and responsibly with others when using technology or shared materials. Why it fits here: The core learning outcome is collaborative responsibility and respectful participation.
\item \textbf{1A-AP-08}: Model daily processes by creating and following algorithms. Why it fits here: Students use rule-based sequences to complete the game or build task.
\item \textbf{1A-AP-11}: Decompose the steps needed to solve a problem into a precise sequence of instructions. Why it fits here: Teacher mediation helps students return to the task sequence after conflict.
\end{itemize}

\section*{NYS / Local Curriculum Alignment}
\begin{itemize}[leftmargin=*]
\item Aligned to the local first-grade IGNITE sequence as a readiness lesson for collaborative computing and maker-space participation.
\item Supports New York-facing SEL and classroom technology participation goals through shared problem solving and material care.
\item This lesson is not programming-heavy; its main CS value is building the classroom culture required for later successful computing tasks.
\end{itemize}

\section*{Lesson Content and Concepts}
\begin{itemize}[leftmargin=*]
\item rules
\item sequence
\item shared resources
\item problem solving
\item collaboration
\end{itemize}

\section*{Evidence / Proof of Instruction}
\begin{itemize}[leftmargin=*]
\item \textbf{Transcript-derived note:} 260211 Grade 1 Gargana Lego Conflict [../../transcript-derived-lessons/260211-grade-1-gargana-lego-conflict.md] The lesson note identifies the class as a LEGO/game context where conflict resolution and task continuation were central.
\item \textbf{Source transcript:} 260211 Grade 1 Gargana Lego Conflict [../../transcripts/260211-grade-1-gargana-lego-conflict.md] The source transcript records redirection, peer conflict management, and attempts to keep the class oriented toward the shared activity.
\item \textbf{Observed teacher moves:} The evidence of instruction comes from norms-setting, conflict mediation, and returning students to the step-by-step activity instead of abandoning the task.
\item \textbf{Likely student proof:} Students completing the game/build, using calmer language, and maintaining access to shared materials are classroom evidence that the lesson landed.
\end{itemize}

\section*{Assessment Evidence}
\begin{itemize}[leftmargin=*]
\item Student stays in the task routine after a conflict or correction.
\item Student can name a rule or repair move that helps the group continue.
\item Student completes their part of the shared game/build structure.
\end{itemize}

\section*{UDL and Inclusion Supports}
\subsection*{Multiple Representation}
\begin{itemize}[leftmargin=*]
\item Keep rules visible with icons and one-sentence language.
\item Model one repair phrase and one help phrase.
\item Use floor/table examples to show what “ready to build” looks like.
\end{itemize}

\subsection*{Multiple Action / Expression}
\begin{itemize}[leftmargin=*]
\item Allow students to point to a rule card instead of speaking immediately.
\item Provide assigned jobs so each student has a clear role.
\item Use quick reset spaces or pause cards for students who become dysregulated.
\end{itemize}

\subsection*{Multiple Engagement}
\begin{itemize}[leftmargin=*]
\item Make the goal concrete and shared so students see why cooperation matters.
\item Use frequent praise for repair, not only for the final product.
\item Keep the task hands-on and short-cycle.
\end{itemize}

\section*{How We Are Already Making It Inclusive}
\begin{itemize}[leftmargin=*]
\item The transcript evidence suggests active teacher repair language and behavioral coaching were already built into the lesson.
\item The task uses tactile materials, which can help students stay engaged while practicing norms.
\item Collaborative play created a natural context for teaching shared-resource expectations.
\end{itemize}

\section*{How to Improve Inclusion Next Time}
\begin{itemize}[leftmargin=*]
\item Add explicit visual sentence stems for conflict repair.
\item Record which roles consistently trigger conflict and redesign those role expectations.
\item Prepare a calmer “restart routine” so students can rejoin without shame.
\item Add a teacher observation checklist to \texttt{artifacts/} for documenting collaboration growth.
\end{itemize}

\section*{Materials and Setup}
\begin{itemize}[leftmargin=*]
\item LEGO or maker materials
\item Game/build prompt
\item Visible rule cards
\item Teacher observation sheet
\end{itemize}

\section*{Teacher Moves / Script / Routines}
\begin{itemize}[leftmargin=*]
\item Preview one rule, one role, and one repair phrase before materials go out.
\item Interrupt quickly when conflict threatens access to the task.
\item Narrate successful repair as evidence of learning, not just behavior compliance.
\item Close with a short reflection on how the group kept the activity going.
\end{itemize}

\section*{Common Errors and Debugging}
\begin{itemize}[leftmargin=*]
\item Students may focus on winning or building faster than working together.
\item Conflict can erase the instructional goal if there is no visible restart routine.
\item Some students may need explicit language frames to re-enter the activity.
\end{itemize}

\section*{Extensions and Simplifications}
\begin{itemize}[leftmargin=*]
\item Extension: have students help generate the next version of the class build rules.
\item Simplification: reduce the number of shared parts and shorten the task cycle.
\item Extension: compare what makes a build/game rule effective versus confusing.
\end{itemize}

\section*{Related Local Materials}
\begin{itemize}[leftmargin=*]
\item No direct local lesson packet linked yet.
\end{itemize}

\section*{Artifacts Checklist}
\begin{itemize}[leftmargin=*]
\item Compiled lesson PDF in \texttt{artifacts/}
\item At least one screenshot, student example, or setup photo
\item One short assessment or observation record
\item Updated standards/source bibliography once the \texttt{.bib} file is revised
\item Any teacher reflection or revision notes after reteaching
\end{itemize}

\section*{Selected Official Sources}
Selected official sources that informed this packet are listed below.
ocite{*}

\bibliographystyle{plain}
\bibliography{lego-game-and-conflict-resolution}
\end{document}
